\subsection{The boundary volume and section growth}
\label{sec:elliptic-boundary-volume}

For the fixed balanced volume on this threefold, the zero potential solves the equation for every positive linear combination of $\alpha,\beta,D$. The variational volume is therefore an explicit cubic polynomial. Its order at the nef boundary records which of the two horizontal coefficients vanish.

\begin{proposition}\label{prop:elliptic-boundary-volume}
On the elliptic threefold, put $\theta=a\alpha+b\beta$ with $a,b\geq0$, and let $\mu_0=\omega^3/(24\pi)$. Then for every $\epsilon>0$,
\begin{equation}\label{eq:elliptic-boundary-volume}
 \mathcal V_{\mu_0}(\theta+\epsilon\omega)
       =24\pi\epsilon(a+\epsilon)(b+\epsilon).
\end{equation}
In particular $\mathcal V_\mu(\theta)=0$ for every smooth positive probability measure $\mu$. If $q$ is the number of positive entries among $a,b$, then the orders of \cref{prop:nef-degeneration-orders} satisfy
\[
 \underline a(\theta)=\overline a(\theta)=3-q.
\]
These orders hold for every smooth positive probability measure and every Hermitian perturbation direction.
\end{proposition}
\begin{proof}
The forms $\alpha,\beta,D$ are semipositive of rank one, their squares vanish, and $\omega=\alpha+\beta+D$ is strictly positive. Thus
\[
 \Omega_\epsilon=\theta+\epsilon\omega
       =(a+\epsilon)\alpha+(b+\epsilon)\beta+\epsilon D>0.
\]
The wedge expansion gives
\[
 \Omega_\epsilon^3
   =6\epsilon(a+\epsilon)(b+\epsilon)\alpha\wedge\beta\wedge D
   =24\pi\epsilon(a+\epsilon)(b+\epsilon)\mu_0.
\]
The zero potential therefore solves its smooth normalized equation. By \cref{thm:gaud-volume-dual-formula}, the equation's scalar is precisely $\mathcal V_{\mu_0}(\Omega_\epsilon)$, proving \eqref{eq:elliptic-boundary-volume}. Nef continuity in \cref{thm:gaud-volume-nef-limit} and smooth-density comparison give $\mathcal V_\mu(\theta)=0$.

If both $a,b$ are positive, the leading term is $24\pi ab\epsilon$. If exactly one is positive, it is $24\pi(a+b)\epsilon^2$. If both are zero, the formula is $24\pi\epsilon^3$. The logarithmic quotients have limits $1,2,3$, respectively. Their independence of the density and perturbing metric is \cref{prop:nef-degeneration-orders}.
\end{proof}

\Needspace{9\baselineskip}
For nonnegative integers $a,b$, the bundle $L=\pi^*\mathcal O_B(a,b)$ has
\[
 h^0(X,L^m)=(am+1)(bm+1).
\]
Indeed the compact-fiber descent and bihomogeneous polynomial argument of \cref{prop:elliptic-selected-section-descent} also apply when one or both integers vanish. Its growth degree is exactly $q$, while the boundary order of its smooth curvature representative is $3-q$. Thus the boundary order records the loss of algebraic growth dimension in this explicit family, even though the unperturbed three-dimensional Monge--Amp\`ere measures vanish.

Every positive perturbation $\Omega_\epsilon$ still has zero lower mass, by \cref{thm:vanishing-lower-mass}, and finite upper mass, since $X$ has BMP. Formula \eqref{eq:elliptic-boundary-volume} concerns one fixed smooth density, and does not assert conservation of the raw mass over all potentials. Whether comparable identifications hold beyond such explicit families is part of \cref{qu:nef-degeneration-exponent}.
